\documentclass[12pt]{article}
\usepackage[utf8]{inputenc}
\usepackage[T1]{fontenc}
\usepackage{hyperref}
\title{Turi Language Resources for K-12 Curriculum Development}
\author{Compiled by Claude Sonnet 5 (claude-sonnet-5) for Andr\'as Kornai}
\date{}
\begin{document}
\maketitle
\section{Introduction}
Turi (ISO 639-3 trd) is a North Munda (Kherwarian) language of India. Estimates run from about 1,500 to 2,000 speakers. The speakers are spread across Chhattisgarh, Jharkhand, Odisha, Bihar, West Bengal and Assam. Many ethnic Turi have shifted to Sadri, Mundari or Odia. A 2024 article notes that children still learn Turi as a home language in northwestern Odisha.
\section{Primary Educational Resources}
No primers, textbooks or school programmes were found in the search results.
\section{Linguistic and Reference Materials}
\begin{itemize}
\item Peterson et al., A Brief Introduction to the Turi Language of Eastern India, Bhasha 3(2), 2024: \url{https://edizionicafoscari.it/it/edizioni4/riviste/bhasha/2024/2/a-brief-introduction-to-the-turi-language-of-easte/}
\item PDF of the same article: \url{https://edizionicafoscari.it/media/pdf/article/bhasha/2024/2/art-10.30687-bhasha-2785-5953-2024-02-005.pdf}
\item Wikipedia, Turi language: \url{https://en.wikipedia.org/wiki/Turi_language}
\item Endangered Languages Project entry: \url{https://endangeredlanguages.com/lang/3507}
\end{itemize}
\section{Religious and Community Resources}
None were found in the search results.
\section{Digital and Audio Resources}
No audio, video or corpus resources were found.
\section{Curriculum Development Context}
The Bhasha article is the main modern linguistic source. It is a sketch for researchers, not teaching material. A curriculum would need new materials developed with Turi communities, and it would have to account for the scattered speaker population.
\end{document}